\documentclass[11pt,a4paper]{report}
\usepackage{exo7hilisit}

\begin{document}

%%%%%%%%%%%%%%%%%%%%%%%%%%%%%%%%%%%%%%%%%%%%%%%%%%%%%%%%%%%%%%%%%%%%%%
%%%%%%%%%%%%%%%%%%%%%%%%%%%%%%%%%%%%%%%%%%%%%%%%%%%%%%%%%%%%%%%%%%%%%%

\entete{Hilisit}{Capacité mathématiques}

\titre{Cours -- \'Equations différentielles}

\bigskip
\bigskip

\begin{quote}
\center
\emph{Les équations différentielles jouent un rôle important en mathématique mais s'appliquent aussi aux autres sciences (physique, biologie,...). L'objectif de ce chapitre est double : comprendre ce nouveau type d'équation et savoir résoudre des équations différentielles simples.}
\end{quote}

\bigskip
\bigskip


\textbf{Sections}
\begin{enumerate}[label=\arabic*.]
  \item \textbf{Primitives}
  
  Thèmes : Rappels sur les primitives qui jouent un rôle important pour la résolution des équations différentielles.
  
  Objectifs : 
  Connaître la définition d'une primitive. 
  Connaître le lien entre deux primitives d'une même fonction. 
  Connaître les formules des primitives usuelles. 
  Savoir déterminer une primitive.
  
  
  \item \textbf{Notion d'équation différentielle}
  
  Thèmes : Introduction et motivation aux équations différentielles.

  Objectifs : 
  Comprendre ce qu'est une équation différentielle. 
  Savoir vérifier qu'une fonction est solution d'une équation différentielle.
  
  \item \textbf{\'Equations différentielles $y'=ay$}
    
  Thèmes : Résolution des équations différentielles $y'=ay$.
  
  Objectifs : 
  Connaître la formule de la solution d'une équation différentielle $y'=ay$.
  Comprendre ce qu'est une condition initiale d'une équation différentielle.
  Comprendre qu'il y a unicité d'une solution lorsqu'on impose une condition initiale.
  Savoir résoudre une équation différentielle $y'=ay$.
  Savoir trouver la solution vérifiant une condition initiale donnée.
  
  \item \textbf{\'Equations différentielles $y'=ay+b$ et $y'=ay+f$}
    
  Thèmes : Résolution des équations différentielles $y'=ay+b$ et $y'=ay+f$
  
  Objectifs : 
  Connaître le vocabulaire : second membre, équation homogène, solution particulière.
  Comprendre la structure d'une solution générale : "solutions de l'équation homogène + solution particulière".
  Savoir trouver une solution particulière à l'aide d'indications.
  Savoir trouver toutes les solutions à partir d'une solution particulière et des solutions de l'équation homogène.
  
\end{enumerate}


\bigskip
\bigskip

\begin{quote}
\center
\emph{Dans ce chapitre on se permettra de noter la fonction $x \mapsto f(x)$ simplement par $f(x)$ pour alléger l'écriture. Par exemple on dira qu'une primitive de $\cos(x)$ est $\sin(x)$ ou bien que $e^x$ est solution de l'équation différentielle $y'=y$.}
\end{quote}

\vfill

\begin{center}
\begin{minipage}{0.8\textwidth}
\center
Auteurs : Arnaud Bodin, Barnabé Croizat et Christine Sacré de l'université de Lille.
Certaines parties sont tirées d'un travail avec Cécile Mammez. Relecture de Pascal Romon.

  \medskip
  
Ce travail a été effectué en 2021-2022 dans le cadre d'un projet Hilisit porté  Unisciel.
\end{minipage}

  \medskip

\raisebox{1ex}{\includegraphics[height=1.8cm]{logo-unisciel}}\qquad\qquad
\includegraphics[height=2.2cm]{logo-ulille}

  \medskip
  
Ce document est diffusé sous la licence \emph{Creative Commons -- BY-NC-SA -- 4.0 FR}.


Sur le site Exo7 vous pouvez récupérer les fichiers sources.

\vspace*{0cm}

\end{center}


\newpage


%%%%%%%%%%%%%%%%%%%%%%%%%%%%%%%%%%%%%%%%%%%%%%%%%%%%%%%%%%%%%%%%%%%%%%
%%%%%%%%%%%%%%%%%%%%%%%%%%%%%%%%%%%%%%%%%%%%%%%%%%%%%%%%%%%%%%%%%%%%%%

\import{./}{equadiff-part1.tex}
\newpage

\import{./}{equadiff-part2.tex}
\newpage

\import{./}{equadiff-part3.tex}
\newpage

\import{./}{equadiff-part4.tex}
\newpage


%%%%%%%%%%%%%%%%%%%%%%%%%%%%%%%%%%%%%%%%%%%%%%%%%%%%%%%%%%%%%%%%%%%%%%
%%%%%%%%%%%%%%%%%%%%%%%%%%%%%%%%%%%%%%%%%%%%%%%%%%%%%%%%%%%%%%%%%%%%%%

\end{document}
